\section*{Bab \ref{ch:b3:rna-regulation} --- RNA Tak Penyandi dan Pengaturan Pascatranskripsi}

\begin{solution}{exo:b3:rna-regulation:1}
miRNA: disandikan di dalam genom sebagai jepit rambut, \qty{22}{nt}, dipotong
\omterm{def:b3:rna-regulation:mirna}{Drosha} lalu \omterm{def:b3:rna-regulation:rnai}{Dicer}, berpasangan tak sempurna lewat benihnya dan menekan
translasi serta kemantapan banyak RNA duta. \omterm{def:b3:rna-regulation:ncrna}{siRNA}: dari RNA beruntai
ganda yang panjang (virus, transgen, percobaan), \qty{21}{nt},
\omterm{def:b3:rna-regulation:rnai}{bergantung-Dicer}, berpasangan sempurna dan mengarahkan \omterm{def:b3:rna-regulation:rnai}{Argonaute} memotong
sasarannya. \omterm{def:b3:rna-regulation:ncrna}{piRNA}: dari transkrip beruntai tunggal \omterm{def:b3:rna-regulation:pirna}{gugus piRNA},
\qtyrange{26}{31}{nt}, tak bergantung \omterm{def:b3:rna-regulation:rnai}{Dicer}, dimuatkan pada protein PIWI di
dalam garis nutfah, memotong transkrip transposon (ping-pong) dan mengarahkan
pembungkaman \omterm{def:b3:chromatin-epigenetics:nucleosome}{kromatin} DNA transposon.
\end{solution}

\begin{solution}{exo:b3:rna-regulation:2}
Pol II mentranskripsi pri-miRNA; \omterm{def:b3:rna-regulation:mirna}{Drosha} (bersama DGCR8) memotong jepit
rambutnya keluar di dalam nukleus; eksportin-5 mengekspor pre-miRNA; \omterm{def:b3:rna-regulation:rnai}{Dicer}
memotong lengkungnya sehingga tersisa duplek \qty{22}{nt}; satu untainya
dimuatkan ke dalam \omterm{def:b3:rna-regulation:rnai}{Argonaute}; benihnya (nukleotida 2--8) berpasangan dengan
tapak UTR 3$'$; \omterm{def:b3:rna-regulation:rnai}{Argonaute} merekrut permesinan deadenilase dan pembuang tudung
serta menghalangi inisiasi --- RNA dutanya lebih jarang ditranslasi dan lebih
cepat meluruh.
\end{solution}

\begin{solution}{exo:b3:rna-regulation:3}
RNA yang dibuat di luar sel oleh polimerase fag memuat cemaran beruntai
ganda (polimerasenya juga menyalin untai lain cetakannya dan berjalan
melampaui ujungnya), sehingga sediaan ``sens'' pun membawa pemicu yang
sebenarnya. Fire dan Mello memurnikan tiap untai, menunjukkan bahwa
masing-masing sendirian hanya sedikit berdaya, lalu menyatukannya dengan
sengaja: RNA beruntai gandanya sedikitnya seratus kali lebih ampuh, dan
beberapa molekul tiap sel sudah memadai.
\end{solution}

\begin{solution}{exo:b3:rna-regulation:4}
Kelaparan besi: IRP mengikat IRE-nya. Feritin: inisiasi translasinya
terhalang, tak ada protein simpanan yang dibuat (\omterm{prop:b3:rna-regulation:translation}{kendali translasi}).
Reseptor transferin: RNA dutanya terlindung dari nukleasenya dan menumpuk,
sehingga lebih banyak reseptor dibuat dan lebih banyak besi diimpor (kendali
kemantapan RNA duta).
\end{solution}

\begin{solution}{exo:b3:rna-regulation:5}
$\gamma = \ln 2/\qty{120}{min} = \qty{0.0058}{min^{-1}}$; $m^{*} =
5/0.0058 \approx 870$ molekul. Peluruhannya dilipatduakan: $m^{*} \approx 430$,
waktu paruhnya $\qty{120}{min}/2 = \qty{60}{min}$.
\end{solution}

\begin{solution}{exo:b3:rna-regulation:6}
Seluruh urutan UTR 3$'$-nya $2\times 10^{7}$ nt; kecocokan kebetulannya $2\times
10^{7}/4^{7} = 2\times 10^{7}/\num{16384} \approx \num{1200}$ tapak.
Sasaran yang sungguhan dibedakan oleh kelestarian tapaknya antarspesies,
konteks yang menguntungkan (tetangga yang kaya AU, tapak yang jauh dari
kodon henti, pasangan tambahan pada ujung 3$'$ miRNA-nya), berbarengnya
ekspresi miRNA dan sasarannya di dalam satu sel, serta percobaan: RNA dutanya
turun ketika miRNA-nya hadir dan tak lagi turun ketika tapaknya dimutasikan.
\end{solution}

\begin{solution}{exo:b3:rna-regulation:7}
(a) XX, \emph{tra} nol: SXL dibuat tetapi TRA tidak, \emph{dsx} mengambil
sambungan baku jantan --- jantan somatik (mandul). (b) XY dengan TRA yang
terus menyala: TRA-2 hadir pada kedua jenis kelamin, sehingga DSX dibuat dalam
bentuk betina --- perkembangan somatik betina (seekor
``pseudobetina'' yang mandul, garis nutfahnya mengikuti petunjuk lain). (c)
Tanpa \emph{dsx}: tidak ada bentuk DSX mana pun, sehingga tak satu pun
program diferensiasi akhir dipaksakan --- seekor interseks dengan ciri
bercampur atau tak lengkap.
\end{solution}

\begin{solution}{exo:b3:rna-regulation:8}
Pada keadaan tunak RNA dutanya, dan karena itu proteinnya, sebanding
dengan waktu paruh RNA dutanya: $\qty{180}{min}/\qty{15}{min} = 12$ kali
lipat lebih banyak protein. Proteinnya adalah isyarat pertumbuhan normal yang
semestinya sekejap, satu denyut sesudah tiap rangsangan; dibuat dua belas kali
lipat dan terus-menerus, ia memacu proliferasi tanpa rangsangan. Dosis dan
penataan waktunya, bukan urutannya, yang menjadikannya onkogenik.
\end{solution}

\begin{solution}{exo:b3:rna-regulation:9}
(a) Induk laboratorium, ayah liar: telurnya tidak membawa \omterm{def:b3:rna-regulation:ncrna}{piRNA} terhadap
elemen P (garis keturunan induknya tak pernah berjumpa dengannya), elemen P
dari ayah bertransposisi bebas di dalam garis nutfah keturunannya --- mandul
(disgenik). (b) Induk liar: telurnya memuat \omterm{def:b3:rna-regulation:ncrna}{piRNA} elemen P, yang
membungkam elemen itu pada keturunannya --- subur. Ketaksetangkupannya adalah
penitipan \omterm{def:b3:rna-regulation:ncrna}{piRNA} oleh induk, yakni pewarisan RNA, bukan pewarisan gen.
\end{solution}

\begin{solution}{exo:b3:rna-regulation:10}
Proteinnya merata-ratakan RNA dutanya sepanjang umurnya sendiri $\tau_{p}$.
RNA duta hidup $\tau_{m} = 1/(\gamma + \gamma' R)$; sepanjang $\tau_{p}$
proteinnya karena itu menyaksikan kira-kira $\tau_{p}/\tau_{m}$ umur RNA duta
yang saling bebas, yang masing-masing merupakan peristiwa lahir--mati acak.
Riak nisbi sebuah rata-rata atas $N$ peristiwa bebas turun sebagai
$1/\sqrt{N}$, sehingga $\tau_{m}$ yang lebih pendek di bawah kendali
\omterm{def:b3:rna-regulation:ncrna}{mikroRNA} --- pada RNA duta rata-rata yang sama, yang menuntut $k$ sebanding
lebih tinggi --- memberi protein yang lebih mulus. Selnya membayar \omterm{def:b3:rna-regulation:ncrna}{mikroRNA}
dan transkripsi tambahan itu demi ketelitian dan kecepatan: kadar yang
ditentukan oleh imbangan sintesis cepat dan peluruhan cepat sekaligus lebih
sedikit berderau dan lebih gesit menyesuaikan diri daripada kadar yang sama
yang ditentukan oleh sintesis lambat dan peluruhan lambat.
\end{solution}

\begin{solution}{exo:b3:rna-regulation:11}
Cacing: RNA polimerase bergantung-RNA membuat \omterm{def:b3:rna-regulation:ncrna}{siRNA} baru dari RNA duta
sasarannya sendiri, sehingga isyaratnya diperbarui selama sasarannya masih
ditranskripsi, bertahan terhadap pengenceran sepanjang pembelahan dan, dengan
adanya saluran RNA beruntai ganda antarsel, menyebar ke jaringan lain serta
ke dalam garis nutfah selama beberapa generasi. Mamalia: tanpa penguatan,
sehingga \omterm{def:b3:rna-regulation:ncrna}{siRNA-nya} terpakai habis dan terencerkan pada tiap pembelahan,
bekerja hanya pada sel yang menerimanya, dan memudar dalam hitungan hari pada
sel yang membelah. Akibatnya: obat \omterm{def:b3:rna-regulation:ncrna}{siRNA} harus diantarkan ke dalam setiap sel
sasaran, bekerja paling baik pada sel yang tidak membelah (hepatosit, karena
itulah obat hati), dan menuntut pemberian berulang.
\end{solution}

\begin{solution}{exo:b3:rna-regulation:12}
(a) RNA yang berfungsi: memusnahkan transkripnya sesudah dibuat (\omterm{def:b3:rna-regulation:ncrna}{siRNA},
oligonukleotida antisens) memberi fenotipe, dan mengekspresikan RNA-nya dari
transgen di tempat lain di dalam genom menyelamatkannya. (b) Transkripsi yang
berfungsi: menyisipkan isyarat poliadenilasi dini yang menghentikan
transkripsi (atau menghapus promotornya) memberi fenotipenya, sedangkan
memusnahkan RNA-nya secara pascatranskripsi tidak, dan transgen ektopik tidak
menyelamatkannya --- tindakan transkripsinya (atau unsur DNA-nya) yang
penting, sebagaimana bagi banyak RNA yang terkait penguat. (c) Derau: tak
satu pun perlakuan mengubah apa pun, RNA-nya kurang dari satu salinan tiap
sel, urutan dan ekspresinya tidak lestari.
\end{solution}

\begin{solution}{pb:b3:rna-regulation:1}
\textbf{1.} $\gamma = \ln 2/40 = \qty{0.0173}{min^{-1}}$; $m^{*} =
3/0.0173 \approx 170$ RNA duta.
\textbf{2.} $\delta_{p} = \ln 2/180 = \qty{0.00385}{min^{-1}}$; $P^{*} =
\beta m^{*}/\delta_{p} = 2\times 173/0.00385 \approx 9.0\times 10^{4}$
protein.
\textbf{3.} Peluruhannya $4\gamma = \qty{0.069}{min^{-1}}$; $m^{*} = 3/0.069
\approx 43$; tetapan waktunya $1/0.069 = \qty{14}{min}$.
\textbf{4.} $\beta = 1$: $P^{*} = 43/0.00385 \approx 1.1\times 10^{4}$;
penekanannya $8$ kali lipat ($4$ dari peluruhan, $2$ dari translasi).
\textbf{5.} Waktu paruh RNA dutanya $\ln 2/(4\gamma) = \qty{10}{min}$.
Proteinnya lalu meluruh dengan waktu paruhnya sendiri, \qty{3}{h}: penguraian
proteinnya yang membatasi sakelar itu.
\textbf{6.} Ya: RNA dutanya hanya turun ke $1/4$ (masih mudah
terdeteksi) sedangkan proteinnya turun ke $1/8$ dan terus turun seiring
protein lamanya berganti; pengamatan awal bahwa proteinnya lenyap ``ketika
mRNA-nya masih ada'' terutama mencerminkan separuh translasi penekanan itu.
\textbf{7.} $10^{6}/250 = 4000$ molekul tiap telur.
\textbf{8.} $4000/550 \approx 7$ tiap sel saat menetas. $4000/2^{n} < 1$
untuk $n \ge 12$ pembelahan sejak zigot --- kira-kira tiga pembelahan
sesudah menetas.
\textbf{9.} Polimerasenya menyalin RNA duta sasarannya menjadi RNA
beruntai ganda baru, yang dipotong menjadi \omterm{def:b3:rna-regulation:ncrna}{siRNA} sekunder: pemicunya
dibangkitkan kembali dari sasarannya sendiri, sehingga ia bertahan terhadap
pengenceran dan menyebar (sebuah saluran meneruskan RNA beruntai ganda
antarsel dan ke dalam garis nutfah). Pengaruhnya memudar karena penguatannya
menuntut sasaran sekaligus pemicu primer; begitu RNA suntikannya habis,
kolam sekundernya terencerkan generasi demi generasi dan garis nutfahnya
mengatur ulang.
\textbf{10.} $5000/2^{n} < 100$: $2^{n} > 50$, $n = 6$ hari ($78$
salinan tersisa).
\textbf{11.} Sel yang tidak membelah tidak mengencerkan \omterm{def:b3:rna-regulation:rnai}{RISC-nya}; batasnya
adalah umur kimia \omterm{def:b3:rna-regulation:ncrna}{siRNA} itu (penguraian oleh nuklease, yang diperlambat
modifikasi kimia untainya) dan rembesan lambat depot endosom yang memasok
sitoplasmanya.
\textbf{12.} $2\times 10^{7}/\num{16384} \approx \num{1200}$ kecocokan
kebetulan. Genetikanya yang memutuskan: kehilangan \emph{lin-14} memberi
fenotipe yang berlawanan dengan kehilangan \emph{lin-4}, alel
\emph{lin-14} yang memperoleh fungsi ternyata berupa penghapusan tepat pada
tapak UTR 3$'$-nya, dan ketujuh tapak itu lestari.
\textbf{13.} Kadarnya $\propto$ waktu paruh: kenaikan $50/10 = 5$ kali lipat.
\textbf{14.} Sintesis feritinnya turun $1/0.05 = 20$ kali lipat.
\textbf{15.} Nisbah impor terhadap simpanan naik $5\times 20 = 100$ kali
lipat dari keadaan kaya besi ke keadaan miskin besi.
\textbf{16.} Gugusnya hanya terbentuk ketika besi sitosol tersedia,
sehingga konformasi proteinnya menjadi pembacaan langsung kepekatan besi,
tanpa reseptor, hormon, atau transkripsi --- seperti sebuah
ribosakelar, sebuah metabolit yang diindera oleh molekul yang dikendalikannya.
\textbf{17.} Feritin ditranslasi tanpa peduli besi: feritin serumnya
sangat tinggi, besi serum dan simpanannya normal (tanpa kelebihan muatan).
Di dalam lensa, yang tak berganti, feritinnya menumpuk bertahun-tahun lalu
mengkristal --- sebuah katarak dini.
\textbf{18.} Jepit rambutnya hanya sebuah tempat mendarat; pengaruhnya
bergantung pada apa yang secara fisik dihalangi protein yang terikat ---
ribosom yang memindai pada ujung 5$'$, sebuah nuklease pada ujung 3$'$.
\textbf{19.} Dua kimia meliputi dua keadaan: gugus besi--belerang
IRP1 juga dimusnahkan oksidan dan oksigen rendah, sehingga IRP1
menanggapi keadaan redoksnya, sedangkan penguraian IRP2 yang bergantung besi
menanggapi besi itu sendiri di sepanjang rentang oksigen fisiologis;
kelewahannya membuat regulonnya tangguh, dan keduanya melaporkan hal yang
berbeda.
\textbf{20.} $12\times 48\times 33\times 2 = \num{38016}$. $1 - (1 -
50/\num{38016})^{50} = 1 - 0.99868^{50} \approx 1 - e^{-0.066} =
0.064$: sekitar \qty{6}{\%}.
\textbf{21.} Isoform yang identik saling mengikat secara homofilik lalu
menolak: cabang satu neuron saling menghindar dan memencar, sedangkan cabang
neuron yang berbeda, yang hampir tak berbagi isoform, boleh bersilangan ---
jati diri pribadi tiap sel.
\textbf{22.} $2^{n}$; $\num{1024}$ untuk $n = 10$. Kenyataannya lebih sedikit:
penyertaannya diputuskan faktor khas tipe sel dan diselaraskan
antarekson, banyak gabungan menggeser kerangka bacanya lalu dimusnahkan
peluruhan yang diperantarai omong kosong, dan ekson-eksonnya tidak saling
bebas.
\textbf{23.} \emph{Sxl} juga mengendalikan \omterm{def:b3:chromatin-epigenetics:x-inactivation}{kompensasi dosis}: pada betina
SXL menghalangi translasi \emph{msl-2}, sehingga kedua kromosom X-nya
tidak ditranskripsi berlebih; tanpa SXL embrio XX merakit kompleks MSL
pada kedua kromosom X dan melipatduakan keluarannya --- kelebihan dosis yang
mematikan bagi setiap gen terpaut-X, bahkan sebelum jenis kelaminnya menjadi
soal.
\textbf{24.} Promotor dininya hanya menyala selama isyarat X:A
hadir; SXL yang dibuatnya memaksa penyambungan produktif transkrip
dari promotor pemeliharaan, yang aktif pada kedua jenis kelamin, sehingga
begitu SXL ada ia terus membuat dirinya sendiri dan isyaratnya tak lagi
dibutuhkan --- logika pembaca-merekrut-penulis yang menopang diri, sama
seperti tanda \omterm{def:b3:chromatin-epigenetics:nucleosome}{kromatin}, di sini dengan sebuah protein dan sebuah sambungan.
\textbf{25.} Protein LIN-14 tertekan $8$ kali lipat; waktu paruh sakelar
RNA dutanya \qty{10}{min}; RNA yang tak diperkuat jatuh di bawah satu molekul
tiap sel sesudah $12$ pembelahan.
\end{solution}
